\section{Introduction}

A striking paradox has emerged in the age of large language models. The same system that passes the bar exam, generates executable code, and proves mathematical theorems can fail catastrophically at tasks a four-year-old performs effortlessly: judging whether a cup is within reach, predicting the outcome of a gentle push, or rerouting around an unexpected obstacle. These are not failures of knowledge---the model likely ``knows'' all the relevant facts. They are failures to translate static knowledge into situated action. An LLM, however capable, is a frozen archive of what has been written. It knows \emph{about} the world, but it does not inhabit it.
% 大语言模型时代出现了一个引人注目的悖论。同一个系统，能通过律师资格考试、生成可执行代码、证明数学定理，却可能在四岁孩童都能轻松完成的任务上灾难性地失败：判断杯子是否在可触及范围内、预测轻轻一推的结果、或绕过意外障碍重新规划路线。这些不是知识的失败——模型很可能"知道"所有相关事实。它们是将静态知识转化为情境化行动的失败。LLM无论多么强大，都只是已书写内容的冻结存档。它"知道"关于世界的事情，但并不栖居于世界之中。

The past two years have witnessed a surge of LLM-based agents that wrap language models with tools, APIs, and memory. Yet despite impressive demos, a pattern of brittleness persists: the web agent fails when a CAPTCHA appears; the robot stalls when an unexpected obstacle blocks its path. This brittleness is not a model problem but a systems problem. The LLM was designed to predict the next token from a static distribution, not to interact with a changing environment. To act, it must be embedded in an external scaffold---an Agent Harness---that supplies observation, execution, and real-time feedback.
% 过去两年间，基于LLM的智能体激增，将语言模型与工具、API和记忆相结合。然而，尽管演示令人印象深刻，脆弱性模式依然存在：验证码出现时网页智能体失败；意外障碍挡住路径时机器人停滞。这种脆弱性不是模型问题而是系统问题。LLM被设计为从静态分布中预测下一个词元，而非与变化的环境交互。要行动，它必须被嵌入一个外部脚手架——智能体框架——以提供观测、执行和实时反馈。

We introduce a four-level hierarchy to make precise what can be learned offline versus what must be observed online. Level 1 (L1), syntactic correctness, concerns well-formedness of symbols. Level 2 (L2), physical plausibility, concerns what events are physically plausible. Level 3 (L3), historical factuality, concerns what has occurred or is documented. These three levels capture regularities learnable from static text. Level 4 (L4), situated truthfulness, is qualitatively different: it concerns what is true \emph{right now}, in this specific environment. L4 is indexical, ephemeral, and requires real-time observation. The gap between L1--L3 and L4 is the \emph{grounding gap}: the structural distance between what a model knows and what an agent must know to act in the present.
% 我们引入一个四层层级来精确区分什么可以离线学习、什么必须在线观测。第一层（L1）语法正确性，关乎符号的形式良好性。第二层（L2）物理合理性，关乎什么事件在物理上是合理的。第三层（L3）历史事实性，关乎什么已经发生或被记录。这三个层次捕捉了可从静态文本中学习的规律性。第四层（L4）情境真实性有质的区别：它关乎"此时此刻"、在这个特定环境中为真的内容。L4是指示性的、瞬时性的，需要实时观测。L1--L3与L4之间的鸿沟就是"落地鸿沟"：模型所知与智能体为在当下行动所必须知道之间的结构性距离。

\textbf{Our central claim is this: L4 knowledge cannot be acquired from static text corpora; frozen parameters contain no L4 information by construction.} Meta-learning is offline generalization over what has been written; it says nothing about online grounding in L4. The Agent Harness bridges this gap: a runtime framework that supplies observation, execution, and feedback. Intelligence emerges from the coupled system of model, harness, and environment, not from the model alone.
% \textbf{我们的核心主张是：L4知识无法从静态文本语料中获得；冻结参数在构造上就不包含L4信息。}元学习是对"已被书写之内容"的离线泛化；它对L4中的在线落地没有说明。智能体框架弥合了这一鸿沟：一个提供观测、执行和反馈的运行时框架。智能涌现自模型、框架和环境耦合的系统，而非模型单独的性质。

We make four contributions. First, we articulate the four-level hierarchy as a diagnostic tool for distinguishing static knowledge from situated state. Second, we formalize grounding using Bayesian inference and distinguish it from inference, simulation, and retrieval. Third, we specify the Agent Harness as an approximate Bayesian decision engine, formalized as \(\mathcal{H} = (\mathcal{S}, \mathcal{O}, \mathcal{A}, \mathcal{B}, \mathcal{T}, \Omega, \Pi)\). Fourth, we derive three architectural principles and validate them empirically on a code-agent benchmark, showing that grounding gains arise from observation fidelity, not model scale. We argue that progress requires shifting investment from bigger models to better harnesses.
% 我们做出四项贡献。第一，阐明四层层级作为区分静态知识与情境状态的诊断工具。第二，使用贝叶斯推理形式化落地，并将其与推理、模拟和检索区分开来。第三，将智能体框架规定为一个近似贝叶斯决策引擎，形式化为\(\mathcal{H} = (\mathcal{S}, \mathcal{O}, \mathcal{A}, \mathcal{B}, \mathcal{T}, \Omega, \Pi)\)。第四，推导三条架构原则，并在代码智能体基准上进行了实证验证，表明落地增益来自观测保真度而非模型规模。我们认为进展需要将投入从更大的模型转向更好的框架。